%! Choosing a Sample: Four Sampling Techniques — Unit 1, Lesson 1.3 — Student Packet Cover Sheet
% PILOT SHAPE (2026-09-06): modeled on AP Statistics lesson 1.4 minus its AP Practice page.
% THREE packet rows — Warm-Up, Guided Notes & Practice, Homework. The homework is the individual
% practice, started in class, and it is SCORED, so its cell is a \blank{}, never NA.
\documentclass[12pt]{article}
\usepackage{saar-article}
\usepackage{saar-boxes}
\usepackage{ltablex}
\keepXColumns
% Measured banner: the band grows to fit the title block, so a title long enough to wrap grows
% the band rather than running out of it (the stats course's \coverbanner; saar-article does not
% carry one, so it is defined here — every cover copies this block verbatim).
\newsavebox{\bannerbox}
\newlength{\bannerht}
\newlength{\coverbannerpad}
\setlength{\coverbannerpad}{0.16in}       % breathing room above and below the text
\newcommand{\coverbanner}[2]{%
  \sbox{\bannerbox}{%
    \begin{minipage}{\dimexpr\paperwidth-1.4in\relax}%
      \centering
      {\color{white}\textbf{\LARGE \CourseName}}\\[4pt]
      {\color{frostmid}\large\textbf{#1}}\\[6pt]
      {\color{white}\textbf{\large #2}}%
    \end{minipage}}%
  \setlength{\bannerht}{\dimexpr\ht\bannerbox+\dp\bannerbox+2\coverbannerpad\relax}%
  \noindent\begin{tikzpicture}[remember picture, overlay]
    \fill[cerulean] (current page.north west)
      rectangle ([yshift=-\bannerht]current page.north east);
    \node[anchor=north, inner sep=0pt]
      at ([yshift=-\coverbannerpad]current page.north) {\usebox{\bannerbox}};
  \end{tikzpicture}\par
  % The text body starts 0.6in below the page edge (geometry top=0.6in), so reserve only what
  % the band overruns that, plus a gap under the band.
  \vspace*{\dimexpr\bannerht-0.6in+0.16in\relax}%
}

\begin{document}

% Full-bleed cerulean banner, sized to the title block.
\coverbanner{Unit 1}{Lesson 1.3 \quad Choosing a Sample: Four Sampling Techniques}

\vspace{0.06in}
% NAMESTRIP: this is the ONLY component that carries the name row.
\namedateperiod
\vspace{0.18in}

% GRADUAL RELEASE: the vocabulary is taught in the guided notes, so the targets name it
% plainly. The standard codes (PS.DC.2a-b; AFDA.DA.2b) stay in the teacher-facing plan.
\begin{learningtargetbox}
\small
\begin{enumerate}[leftmargin=1.5em, itemsep=5pt]
  \item \ldots{} name the \textbf{sampling frame} and say whether \textbf{chance} or
        \textbf{convenience} chose a sample --- by asking \emph{who could never have been picked}.
  \item \ldots{} carry out a \textbf{simple random sample} from generator digits and a
        \textbf{systematic sample} with interval $k = N \div n$.
  \item \ldots{} compute a \textbf{stratified} allocation in proportion, and count what a
        \textbf{cluster sample} reaches.
  \item \ldots{} tell a \textbf{stratum} from a \textbf{cluster} --- some from every group
        versus everyone from a few --- and choose and \emph{justify} the technique a context
        calls for.
\end{enumerate}
\end{learningtargetbox}

\vspace{0.14in}

\begin{tocbox}
\small
\renewcommand{\arraystretch}{1.5}
\begin{tabularx}{\linewidth}{@{} c l X r @{}}
\rowcolor{frost}
\textbf{\#} & \textbf{Component} & \textbf{Description} & \textbf{Score} \\
\midrule
1 & Warm-Up & Riverbend's $900$ students: a percent, a share of $60$, and an interval down
             a list & \blank{1.2cm} \\
2 & Guided Notes \& Practice & Who chooses the sample; the four techniques; Maya's homeroom
             draw --- then Cedar Ridge Apartments worked together & \blank{1.2cm} \\
% Row 3 is the lesson's GRADED practice, started in class. Packet pages are the default; a
% Desmos activity may replace them on a given day, decided at Close & Assign. The row and its
% score cell never change.
3 & Homework & Harbor Point pool, done properly: run four plans and choose the one the board
             needs --- \textbf{scored, started in class, due the first class after two study
             halls} & \blank{1.2cm} \\
\midrule
  & \multicolumn{2}{r}{\textbf{Total}} & \blank{1.2cm} \\
\end{tabularx}
\end{tocbox}

\vspace{0.14in}

% Keep in Mind carries the lesson's CONTENT — the rule, the test, the trap — never its process.
\begin{remindbox}
\small
In a \textbf{probability sample}, chance chooses --- ask \emph{who could never have been
picked?} \textbf{Simple random}: number the frame and draw, skipping repeats and out-of-range
values. \textbf{Systematic}: a random start, then every $k$th, $k = N \div n$.
\textbf{Stratified}: some from \emph{every} group, in proportion --- for groups that are
\emph{alike inside}. \textbf{Cluster}: everyone from a \emph{few} groups --- only for groups
that are each a \emph{small mix of the whole}. \textbf{Watch out:} ``chance chose'' is not
enough --- draw a few groups that are alike inside and you can miss a whole grade.
\end{remindbox}

\end{document}
